% 2012 年江西卷文科第 10 题：几何图形与 S(t) 的四个图象选项。
% 来源 PDF：普通高考/2012/2012江西文.pdf，第 1 页（印面 749）右栏。
% 几何：O(0,0)，A(2,0)，B 在 30° 方向且 OB=1，圆弧以 A 为圆心、AB≈1.24 为半径交 OA 延长线于 C。
% 正确图象 (B)：0≤t≤1 为抛物线，之后为直线（斜率更大），到 C 后水平。
\begin{tikzpicture}[scale=0.8, semithick, >=Stealth, font=\small]
  % ---------- 几何图形 ----------
  \begin{scope}
    \draw (0,0) node[below left] {$O$} -- (2,0) node[below] {$A$};
    \draw (0.866,0.5) node[above left] {$B$} -- (2,0);
    \draw (0,0) -- (0.866,0.5);
    \draw[thick] (0.866,0.5) arc[start angle=156.2, end angle=0, radius=1.239];
    \draw (2,0) -- (3.24,0) node[below] {$C$};
    \node[above] at (2.0,1.4) {$D$};
    \fill (0.866,0.5) circle (1.5pt);
    \fill (2,0) circle (1.5pt);
    \node[below] at (1.6,-1.35) {};
  \end{scope}
  % ---------- (A) ----------
  \begin{scope}[shift={(0.2,-3.9)}]
    \draw[->] (0,0) -- (3.2,0);
    \draw[->] (0,0) -- (0,2.4);
    \draw[smooth] (0,0) .. controls (0.6,0.05) and (1.2,0.9) .. (1.9,2.0) .. controls (2.2,2.25) and (2.5,2.3) .. (3.0,2.3);
    \node[below left] at (0,0) {$O$};
    \node at (1.6,-0.55) {(A)};
  \end{scope}
  % ---------- (B) ----------
  \begin{scope}[shift={(4.3,-3.9)}]
    \draw[->] (0,0) -- (3.2,0);
    \draw[->] (0,0) -- (0,2.4);
    \draw[smooth, domain=0:1] plot (\x, {0.5*\x*\x});
    \draw (1,0.5) -- (2.2,2.05) -- (3.0,2.05);
    \node[below left] at (0,0) {$O$};
    \node at (1.6,-0.55) {(B)};
  \end{scope}
  % ---------- (C) ----------
  \begin{scope}[shift={(0.2,-7.1)}]
    \draw[->] (0,0) -- (3.2,0);
    \draw[->] (0,0) -- (0,2.4);
    \draw[smooth] (0.3,0.05) .. controls (0.9,0.1) and (1.2,1.9) .. (1.9,2.25) .. controls (2.3,2.35) and (2.6,2.3) .. (3.0,2.3);
    \node[below left] at (0,0) {$O$};
    \node at (1.6,-0.55) {(C)};
  \end{scope}
  % ---------- (D) ----------
  \begin{scope}[shift={(4.3,-7.1)}]
    \draw[->] (0,0) -- (3.2,0);
    \draw[->] (0,0) -- (0,2.4);
    \draw (0.9,0) -- (0.9,2.2) -- (2.2,2.2);
    \node[below left] at (0,0) {$O$};
    \node at (1.6,-0.55) {(D)};
  \end{scope}
\end{tikzpicture}
